\documentclass[11pt,letterpaper]{article}
\usepackage[margin=0.9in]{geometry}
\usepackage{booktabs}
\usepackage{array}
\usepackage{enumitem}
\usepackage{xcolor}
\usepackage{titlesec}
\usepackage{fancyhdr}
\usepackage{hyperref}
\usepackage{tabularx}
\usepackage{tcolorbox}

% Colors
\definecolor{headerblue}{RGB}{47,68,117}
\definecolor{accentblue}{RGB}{70,100,160}
\definecolor{darkgreen}{RGB}{40,100,60}
\definecolor{examplegray}{RGB}{245,245,248}

% Header/Footer
\pagestyle{fancy}
\fancyhf{}
\fancyhead[L]{\small\textcolor{headerblue}{MUS 262 --- Listening Response Framework}}
\fancyhead[R]{\small\textcolor{headerblue}{Reference Sheet}}
\fancyfoot[C]{\thepage}
\renewcommand{\headrulewidth}{0.5pt}

% Section formatting
\titleformat{\section}{\Large\bfseries\color{headerblue}}{}{0em}{}[\vspace{-0.5em}\rule{\textwidth}{0.5pt}]
\titleformat{\subsection}{\large\bfseries\color{accentblue}}{}{0em}{}
\titleformat{\subsubsection}{\normalsize\bfseries}{}{0em}{}

\hypersetup{colorlinks=true, linkcolor=headerblue, urlcolor=accentblue}

\setlength{\parindent}{0pt}
\setlength{\parskip}{0.5em}

\tcbset{
  examplebox/.style={
    colback=examplegray,
    colframe=accentblue,
    boxrule=0.5pt,
    arc=2pt,
    left=6pt, right=6pt, top=4pt, bottom=4pt,
    fonttitle=\bfseries\small,
  }
}

\begin{document}

\begin{center}
{\Huge\bfseries\textcolor{headerblue}{Listening Response Framework}}\\[0.3em]
{\Large What to Say About What You Hear}\\[0.5em]
{\large MUS/AAS 262 --- Jazz History: Many Sounds, Many Voices}
\end{center}

\vspace{0.5em}

\section{The Paragraph Structure}

Your listening response is essentially: \textbf{What do I hear? $\rightarrow$ What does it mean?}

Write 5--6 sentences hitting these categories in order:

\begin{enumerate}[leftmargin=1.5em, itemsep=0.3em]
  \item \textbf{Instrumentation \& Ensemble} --- What instruments? How many players?
  \item \textbf{Tempo \& Rhythm} --- How fast? What's the rhythmic feel?
  \item \textbf{Timbre \& Sound Quality} --- What does it \emph{sound like}? (see below)
  \item \textbf{Melody \& Improvisation} --- Is the melody stated or improvised? How?
  \item \textbf{Form \& Texture} --- What's the structure? Dense or sparse?
  \item \textbf{Connection} --- What era/concept from class does this connect to?
\end{enumerate}

\begin{tcolorbox}[examplebox, title={Template sentence}]
``This piece features \underline{\hspace{2cm}} [instruments], performing at a \underline{\hspace{2cm}} [tempo] with a \underline{\hspace{2cm}} [rhythmic feel]. The soloist's timbre is \underline{\hspace{2cm}}, with \underline{\hspace{2cm}} [vibrato/articulation]. The form is \underline{\hspace{2cm}}, and a notable feature is \underline{\hspace{2cm}}, which connects to \underline{\hspace{2cm}} [course concept].''
\end{tcolorbox}

\section{The Six Categories: What to Listen For \& What to Say}

\subsection{1. Instrumentation \& Ensemble}

\textbf{What you're identifying:} Which instruments are playing and how the group is organized.

\renewcommand{\arraystretch}{1.3}
\begin{small}
\begin{tabularx}{\textwidth}{>{\bfseries}p{3cm} X}
\toprule
\textbf{If you hear\ldots} & \textbf{You can say\ldots} \\
\midrule
One instrument alone & ``solo piano,'' ``unaccompanied trumpet'' \\
A few instruments (3--6) & ``small combo/ensemble,'' ``quartet,'' ``quintet'' \\
Many instruments, sections trading off & ``big band,'' ``full orchestra,'' ``horn sections'' \\
Trumpet/cornet, clarinet, trombone & ``New Orleans front line'' (= Early Jazz) \\
Saxophone sections & ``reed section'' (= Swing era or later) \\
Rhythm section only & ``piano trio,'' ``rhythm section'' (piano, bass, drums) \\
Voice as lead & ``vocal performance with [accompaniment]'' \\
\bottomrule
\end{tabularx}
\end{small}

\textbf{Also note:} Who is the soloist? Is it one person or does the solo pass between players?

\subsection{2. Tempo \& Rhythm}

\textbf{What you're identifying:} How fast it is and how the beat feels.

\begin{small}
\begin{tabularx}{\textwidth}{>{\bfseries}p{3cm} X}
\toprule
\textbf{What to describe} & \textbf{Words you can use} \\
\midrule
Speed & slow / ballad tempo, moderate / medium tempo, up-tempo, very fast / breakneck \\
Swing feel & ``swinging'' (long-short pattern), ``straight'' (even eighth notes), ``stiff'' (ragtime-like) \\
Beat emphasis & ``emphasizes beats 2 and 4'' (jazz standard), ``four-on-the-floor'' (every beat even) \\
Rhythmic energy & ``driving,'' ``relaxed / laid-back,'' ``pushing forward,'' ``behind the beat'' \\
Syncopation & ``syncopated'' (accents on unexpected beats), ``rhythmically complex'' \\
Groove & ``danceable,'' ``bouncy,'' ``relentless,'' ``floating'' \\
\bottomrule
\end{tabularx}
\end{small}

\begin{tcolorbox}[examplebox, title={Examples}]
``The tempo is breakneck, with the rhythm section pushing forward relentlessly.''\\
``A relaxed, medium-tempo swing feel with the drummer emphasizing beats 2 and 4 on the hi-hat.''\\
``The rhythm feels stiff and march-like, characteristic of early ragtime influence.''
\end{tcolorbox}

\subsection{3. Timbre (What Does It \emph{Sound} Like?)}

\textbf{Timbre} (pronounced ``TAM-ber'') = the \textbf{color or quality} of a sound. It's \emph{why} a trumpet sounds different from a saxophone even when they play the same note at the same volume. Think of it as the ``texture'' of the sound itself.

\textbf{Simple way to think about it:} If someone asked ``describe this instrument's voice,'' timbre is your answer.

\begin{small}
\begin{tabularx}{\textwidth}{>{\bfseries}p{3cm} X}
\toprule
\textbf{Timbre dimension} & \textbf{Words you can use} \\
\midrule
Brightness & bright / brilliant / piercing \textbf{vs.} dark / warm / mellow / rich \\
Weight & heavy / thick / full \textbf{vs.} light / thin / airy \\
Texture & smooth / clean / pure \textbf{vs.} rough / gritty / raspy / buzzy / growling \\
Breathiness & breathy / airy \textbf{vs.} focused / clear / sharp \\
Resonance & resonant / ringing / booming \textbf{vs.} dry / muted / dampened \\
Mood descriptors & yearning, plaintive, joyful, aggressive, tender, strident, wistful, sober \\
\bottomrule
\end{tabularx}
\end{small}

\begin{tcolorbox}[examplebox, title={Timbre examples from your listening}]
\textbf{Armstrong's trumpet} = bright, strident, piercing\\
\textbf{Lester Young's sax} = light, airy, cool, smooth\\
\textbf{Coleman Hawkins' sax} = warm, full, lush, heavy vibrato\\
\textbf{Billie Holiday's voice} = fragile, dark, horn-like, yearning\\
\textbf{Bessie Smith's voice} = dark, powerful, heavy vibrato\\
\textbf{Muted trumpet} = dampened, nasal, buzzy (vs. open trumpet = bright, ringing)
\end{tcolorbox}

\subsection{4. Vibrato \& Articulation}

These are two separate things. Both describe \emph{how} notes are played.

\subsubsection{Vibrato}
\textbf{What it is:} A slight, rapid wavering of pitch on a sustained note. Think of a singer holding a long note and it ``wobbles'' --- that's vibrato.

\begin{small}
\begin{tabularx}{\textwidth}{>{\bfseries}p{3cm} X}
\toprule
\textbf{What you hear} & \textbf{What to say} \\
\midrule
Lots of wobble & ``extensive/heavy/wide vibrato'' \\
Subtle wobble & ``gentle/controlled/narrow vibrato'' \\
Almost none & ``minimal vibrato,'' ``straight tone'' \\
Only at end of notes & ``terminal vibrato'' (this is Armstrong's innovation) \\
Throughout the note & ``continuous vibrato'' (Hawkins, Bessie Smith) \\
\bottomrule
\end{tabularx}
\end{small}

\subsubsection{Articulation}
\textbf{What it is:} How each individual note begins and ends --- is it punched, slurred, bitten off?

\begin{small}
\begin{tabularx}{\textwidth}{>{\bfseries}p{3cm} X}
\toprule
\textbf{What you hear} & \textbf{What to say} \\
\midrule
Notes flow together & ``legato'' (smooth, connected) \\
Notes are separated/punchy & ``staccato'' (short, detached, clipped) \\
Sharp attack on each note & ``crisp/hard articulation,'' ``tongued'' \\
Soft starts & ``soft attack,'' ``gentle onset'' \\
Notes slide between pitches & ``glissando,'' ``sliding,'' ``scooping into notes,'' ``bending'' \\
Notes cut off abruptly & ``clipped,'' ``bitten off'' \\
Mix of smooth and punchy & ``varied articulation'' \\
\bottomrule
\end{tabularx}
\end{small}

\begin{tcolorbox}[examplebox, title={Key distinctions from class}]
\textbf{Parker (bebop)}: clear, crisp articulation even at extreme speed --- every note distinguishable\\
\textbf{Billie Holiday}: extreme legato --- words and notes spill into each other\\
\textbf{Bessie Smith}: bending/sliding into notes (blue notes) --- instrumentalists copied this\\
\textbf{Lester Young}: smooth legato lines, behind the beat, relaxed\\
\textbf{Hawkins}: smooth but active, arpeggiated runs up and down
\end{tcolorbox}

\subsection{5. Melody, Improvisation \& Harmony}

\begin{small}
\begin{tabularx}{\textwidth}{>{\bfseries}p{3cm} X}
\toprule
\textbf{What to describe} & \textbf{Words you can use} \\
\midrule
Melody stated clearly & ``the melody/head is stated clearly,'' ``recognizable tune'' \\
Melody abandoned & ``the soloist departs from the melody,'' ``improvises freely over the changes'' \\
Melodic shape & ``angular'' (big jumps), ``stepwise'' (notes close together), ``lyrical,'' ``singable'' \\
Improvisation style & ``collective improvisation'' (everyone at once), ``solo improvisation'' (one person) \\
Phrasing & ``long flowing lines,'' ``short punchy phrases,'' ``call and response,'' ``conversational'' \\
Harmonic complexity & ``simple blues harmony,'' ``complex chord changes,'' ``playing the changes'' \\
Register & ``high register'' (Dizzy), ``middle register,'' ``low register,'' ``altissimo'' (extreme high) \\
\bottomrule
\end{tabularx}
\end{small}

\begin{tcolorbox}[examplebox, title={Vertical vs. Horizontal (a key distinction)}]
\textbf{Vertical} (Hawkins) = playing up and down the notes \emph{within} each chord (arpeggios). You hear the chord changes clearly because the soloist outlines them.\\[0.3em]
\textbf{Horizontal} (Young) = playing melodic lines that flow \emph{across} chord changes. Sounds more like storytelling; you hear a melody, not individual chords.
\end{tcolorbox}

\subsection{6. Form, Texture \& Special Features}

\begin{small}
\begin{tabularx}{\textwidth}{>{\bfseries}p{3cm} X}
\toprule
\textbf{What to describe} & \textbf{Words you can use} \\
\midrule
Form/structure & ``12-bar blues (A,A,B),'' ``32-bar AABA,'' ``multi-strain (ragtime),'' ``contrafact'' \\
Texture (density) & ``dense/thick'' (many instruments), ``sparse/open'' (few instruments, space) \\
Dynamics & ``loud/forte,'' ``soft/piano,'' ``builds in intensity,'' ``dynamic contrast'' \\
Call and response & ``call-and-response between [X] and [Y]'' \\
Special features & cadenza, spoken intro, scat singing, muted brass, drum feature, walking bass \\
Arrangement & ``tightly arranged,'' ``head arrangement,'' ``polished,'' ``loose/spontaneous'' \\
\bottomrule
\end{tabularx}
\end{small}

\section{Putting It All Together: Example Paragraphs}

\begin{tcolorbox}[examplebox, title={Example: Louis Armstrong --- ``West End Blues'' (1928)}]
This piece opens with an \textbf{unaccompanied trumpet cadenza} --- Armstrong playing alone with no rhythm section, demonstrating virtuosic technique. The tempo settles into a \textbf{slow, relaxed 12-bar blues}. Armstrong's trumpet timbre is \textbf{bright and strident}, with \textbf{vibrato applied at the end of sustained notes} rather than throughout --- a technique he pioneered. His articulation alternates between \textbf{smooth legato phrases and crisp, punched notes}. The texture shifts between \textbf{solo trumpet and the full ensemble} (clarinet, trombone, piano), with the band providing a \textbf{sparse} backdrop during solos. This piece exemplifies Armstrong as the ``First Great Soloist'' (Schuller), transforming jazz from collective ensemble music into a vehicle for individual expression.
\end{tcolorbox}

\begin{tcolorbox}[examplebox, title={Example: Coleman Hawkins --- ``Body and Soul'' (1939)}]
This is a \textbf{tenor saxophone} with rhythm section (piano, bass, drums) performing at a \textbf{slow ballad tempo}. Hawkins's timbre is \textbf{warm, full, and lush}, with \textbf{extensive vibrato throughout} his sustained notes. He \textbf{barely states the original melody}, instead improvising almost entirely over the 32-bar AABA chord changes. His approach is \textbf{vertical} --- you can hear him running \textbf{arpeggios up and down the chord tones}, reaching into the \textbf{altissimo register}. The texture is \textbf{sparse}, with the rhythm section staying out of the way. This recording demonstrates Hawkins as the ``Father of the Jazz Tenor Saxophone'' and his chordal/vertical approach that contrasts directly with Lester Young's horizontal style.
\end{tcolorbox}

\begin{tcolorbox}[examplebox, title={Example: Charlie Parker --- ``Ko Ko'' (1945)}]
A small \textbf{bebop combo} --- alto saxophone (Parker), trumpet, piano (Miles Davis), bass, and drums (Max Roach). The tempo is \textbf{extremely fast, almost breakneck}. Parker's timbre is \textbf{bright and focused}, and his articulation is remarkably \textbf{crisp and clear even at this speed} --- every note is distinguishable. The piece opens with a \textbf{pre-composed introduction} before Parker launches into an \textbf{extended improvised solo} over ``I Got Rhythm'' chord changes (a contrafact). His melodic lines are \textbf{angular}, with large intervallic leaps and \textbf{syncopated, unpredictable phrasing}. Max Roach drives the rhythm from the \textbf{ride cymbal} rather than the bass drum, characteristic of bebop drumming. This is bebop as ``art music, not dance music'' --- the complexity and speed demand focused listening, not dancing.
\end{tcolorbox}

\section{Cheat Sheet: Quick Vocabulary by Category}

\begin{small}
\begin{tabularx}{\textwidth}{>{\bfseries\color{headerblue}}p{2.5cm} X}
\toprule
\textbf{Category} & \textbf{Go-to words} \\
\midrule
Tempo & slow/ballad, moderate, up-tempo, breakneck/blazing \\
Rhythm & swinging, straight, stiff, syncopated, driving, laid-back, behind the beat \\
Timbre & bright vs. dark, warm vs. piercing, breathy vs. focused, smooth vs. gritty, lush, fragile \\
Vibrato & extensive/heavy, gentle/subtle, minimal/none, terminal (end of note only) \\
Articulation & legato (smooth) vs. staccato (punchy), crisp, sliding/bending, clipped \\
Melody & angular vs. stepwise vs. lyrical, stated vs. improvised, long lines vs. short phrases \\
Improv & collective vs. solo, vertical (chordal) vs. horizontal (melodic), ``playing the changes'' \\
Texture & dense vs. sparse, full band vs. stripped down, section playing vs. solo \\
Dynamics & quiet/soft, loud/powerful, builds in intensity, dynamic contrast \\
Form & 12-bar blues, 32-bar AABA, contrafact, head arrangement, multi-strain \\
\bottomrule
\end{tabularx}
\end{small}

\end{document}
